\section{Problem Formulation}
\label{sec:problem}

We study a multi-turn exchange in which one scalar property of a language model's behavior is
tracked over time and a controller may issue one intervention per turn under a fixed budget.
An independent judge scores every reply on that property, and that judged score is the only
readout this paper reports; a controller may act on its own in-loop signal, but no reported number
comes from it. Every result is measured against an equal-cost fixed schedule, which spends the
same budget at turn positions chosen without reading the conversation. Our four task families
share one shape, a text property pushed by a reminder or a rewrite, so the results below speak to
those settings and the preconditions we name, not to multi-turn control in general.

\subsection*{Formalization Entrance}

Consider a reader who asks for a paragraph to be simplified one step at a time. Where the
paragraph sits now decides whether the next instruction should push it down a full level, down
half a level, or leave it alone, so the instruction sequence cannot be fixed in advance. Yet an
instruction that moves the text down a level moves it down that level whether it arrives at the
second request or the fifth. The first half of this picture makes a state worth tracking; the
second makes the timing of an action worth less than it appears.

\textbf{Assumption 1 (Readout Persistence).} A window of past readouts carries information about
the next readout that the current turn alone does not. This is a tendency over a task family,
not an identity on every trajectory.

\textbf{Assumption 2 (Action Identifiability).} The intervention varies independently enough of
the state that its effect on the next state can be estimated separately from the state's own
evolution. We state it as an approximation, and we state it rather than dropping it because
Section~\ref{sec:experiments} reports that it fails on every column we can test.

\textbf{Assumption 3 (Identification Predicts Control).} If withholding the action does not
change how well the fitted operator predicts the trajectory, then no controller built on that
operator beats an equal-cost fixed schedule. This is a tendency, and we check one of its two
directions: a zero licenses the negative inference, while the non-zero side stays untested for
the reason given in Section~\ref{sec:conclusion}.

\subsection*{Definitions}

\textbf{Readout and Intervention.} The readout $y_t \in \mathbb{R}$ is the judged score at turn
$t$, and the intervention $v_t$ is what the controller issues there, binary for the reminder
settings and four-valued for the rewriting setting. On the trajectory-tracking tasks the control
is instead $u_t = \phi(r - y_t)$, the error against a per-trajectory target $r$ held fixed for
the whole trajectory, with $\phi$ either the identity or the map appending the error's magnitude
and sign.

\textbf{Delay-Embedded State.} A single readout is not a state. Formalizing Assumption 1, we
collect $L_y$ past readouts and $L_u$ past controls into one vector,
%
\begin{equation}
  z_t = \bigl[\,y_t,\ y_{t-1},\ \dots,\ y_{t-(L_y-1)},\ u_t,\ u_{t-1},\ \dots,\ u_{t-(L_u-1)}\,\bigr],
  \label{eq:delay-state}
\end{equation}
%
where $L_y$ counts readout terms and $L_u$ control terms. Setting $(L_y, L_u) = (1,0)$ gives the
memoryless baseline, and every deeper setting adds one term of each. A turn contributes a row
only when its lagged terms all exist and turn $t+1$ is observed. The row count therefore falls
as the window deepens, and Section~\ref{sec:experiments} shows that this bookkeeping carries a
share of what depth appears to buy.

\textbf{Reported Quantity.} We score a model by how much of a trivial baseline's $H$-step rollout
error it removes,
%
\begin{equation}
  \mathrm{Skill}_H = 1 - \frac{\mathrm{MSE}_H(\text{model})}{\mathrm{MSE}_H(\text{null}^\star)},
  \qquad
  \text{null}^\star = \arg\min_{n \in \{\text{const},\, \text{turn-mean},\, \text{stateless}\}} \mathrm{MSE}_H(n),
  \label{eq:skill}
\end{equation}
%
where $H$ is the rollout horizon and $\text{null}^\star$ is the best of three trivial baselines
fitted on the training split and detailed in Section~\ref{sec:experiments}. Each receives the turn
index, since a baseline denied that index can be beaten by a model that has learned only the ramp.

Equation~\eqref{eq:skill} returns no number in two cases, and both carry content. It rejects
$H < 2$, which keeps one-step prediction error out of every headline. It also leaves
$\mathrm{Skill}_H$ undefined, rather than zero, when $\mathrm{MSE}_H(\text{null}^\star)$ sits at
floor: the readout then has no range, so the operator has nothing to be skillful about, and a
zero would read as a failure of the model instead of an absence in the data. Our negative-control
column is this case. Because each column normalizes against its own baseline, $\mathrm{Skill}_H$
does not compare across columns, and no table here carries a row that averages or ranks over them.

Section~\ref{sec:method} formalizes Assumption 2 in the operator's action channel and
Assumption 3 in the diagnostics built on the fitted operator.
